% Part: computability
% Chapter: recursive-functions
% Section: pr-functions-computable

\documentclass[../../../include/open-logic-section]{subfiles}

\begin{document}

\olfileid{cmp}{rec}{cmp}
\olsection{الدوال العودية البدائية قابلة للحساب}

لتكن الدالة $h$ معرّفة بالعودية البدائية كما يأتي:
\begin{eqnarray*}
h(\vec x, 0)     & = & f(\vec x) \\
h(\vec x, y + 1) & = & g(\vec x, y, h(\vec x, y))
\end{eqnarray*}
ولنفرض قابلية $f$ و$g$ للحساب. والرمز $\vec x$ عندنا اختصار
لـ$x_0$، …،~$x_{k-1}$. أما $h(\vec x,0)$ فيتأتى حسابها
لأنها عين $f(\vec x)$، وهي قابلة للحساب بالفرض. ثم يتأتى
حساب $h(\vec x,1)$؛ إذ $1=0+1$، فتكون $h(\vec x,1)$
\begin{align*}
 h(\vec x, 1) & = g(\vec x, 0, h(\vec x, 0)) =  g(\vec x, 0, f(\vec x)).
\intertext{ويمكننا المواصلة على هذا النحو وحساب}
h(\vec x, 2) & = g(\vec x, 1, h(\vec x, 1)) = g(\vec x, 1, g(\vec x, 0, f(\vec x)))\\
h(\vec x, 3) & = g(\vec x, 2, h(\vec x, 2)) = g(\vec x, 2, g(\vec x, 1, g(\vec x, 0, f(\vec x))))\\
h(\vec x, 4) & = g(\vec x, 3, h(\vec x, 3)) = g(\vec x, 3, g(\vec x, 2, g(\vec x, 1, g(\vec x, 0, f(\vec x)))))\\
& \vdots
\end{align*}
فإذا طلبنا $h(\vec x,y)$ لأي حجة، حسبنا على التعاقب $h(\vec x,0)$ و
$h(\vec x,1)$ و…، حتى نبلغ $h(\vec x,y)$.

فقد بان أن التعريف بالعودية البدائية يحفظ قابلية الحساب متى كانت
الدالتان $f$ و$g$ قابلتين للحساب. وكذلك التركيب يحفظ هذه
القابلية إذا تحققت في $f$ وفي الدوال $g_i$.

وحيث إن الدوال الأساسية $\Zero$ و$\Succ$ و$\Proj{n}{i}$ قابلة
للحساب، وإن التركيب والعودية البدائية لا يخرجان بنا من
الدوال القابلة للحساب، لزم أن تكون كل دالة عودية بدائية قابلة للحساب.

\end{document}
